\section{Background}
\subsection{Context-Free Grammar}

A \emph{context-free grammar} (CFG) consists of terminals, non-terminals, a start symbol, and production rules~\cite{hopcroft2001introduction, Aho_Aho_2007}. It can be formally defined as a quadruple $\cfg = (V, \Sigma, \Pi, S)$, where $V$ is a finite set of non-terminal symbols, $\Sigma$ is the set of terminals, $\Sigma \cap V = \emptyset$, $\Pi$ is a finite set of production rules, $\Pi \subseteq V \times (V \cup \Sigma)^\ast$, and $S \in V$ is the start symbol of the grammar.
Elements of $(V \cup \Sigma)^*$ are known as sentential forms.

The language generated by $\cfg$, denoted $\languageOf{\cfg}$, comprises all strings derivable from $S$ using the rules in $\Pi$: 
$
\languageOf{\cfg} = \{ \sigma \in \Sigma^\ast \mid S \stackrel{*}{\rightarrow} \sigma \}
$. For $\alpha,\beta \in (V\cup \Sigma)^*$, we say $\alpha$ directly derives $\beta$ in one step as $\alpha {\rightarrow} \beta$, and define $\alpha \stackrel{*}{\rightarrow} \beta$ as $\alpha$ deriving $\beta$ in zero or more steps if there exists a finite sequence of $\gamma_0,\dots,\gamma_n \in (V \cup \Sigma)^*$ where $n\geq 0$, such that $\alpha = \gamma_0 \rightarrow \gamma_1 \rightarrow \dots \rightarrow \gamma_n = \beta$.


\subsection{Backus-Naur Form}
\label{sec:bnf_background}
Backus–Naur Form (BNF) is a notation used to define the grammar of formal languages, typically CFG~\cite{Chomsky_1956,backus1959syntax,backus1960report}.  
In this paper, LLM generates grammars in BNF. 
% In BNF, production rules are defined  

A context-free grammar $\cfg = (V, \Sigma, \Pi, S)$ is given in BNF notation as a list of grouped production rules, where each rule group is written:

$$\texttt{<$v_i$>}\vcentcolon\vcentcolon=\alpha_1 | \alpha_2 | \ldots$$

where $v_i$ enclosed between the pair ``\texttt{< >}" is a nonterminal symbol in $V$, and $\alpha_1, \alpha_2 \ldots \in (\Sigma \cup V)^*$ is the list of sentential forms that can be derived in one step from $v_i$. This represents the set of all production rules with $v_i$ on the left-hand side, i.e., $v_i \rightarrow \alpha_1, v_i \rightarrow \alpha_2, \ldots$. 

We extend the definition of a CFG to define a grammar in BNF to be $\bnf = (V, \Sigma, \Pi, S, R)$, where $V, \Sigma, \Pi$, and $S$ are defined as before. $R$ is a set of sets of production rules, denoted $\{r_1, \ldots r_n\}$. Each $r_i \in R$ is a set of production rules where the left-hand side of all rules is the $i^{th}$ non-terminal symbol $v_i$, i.e., $r_i$ is $(v_i \times (\Sigma \cup V)^*) \cap \Pi$. 
Since each non-terminal symbol has a corresponding rule set, $n=|V|=|R|$. The language of the BNF grammar $\bnf$ is denoted by $\languageOf{\bnf}$ and is equal to the language of the original CFG, $\languageOf{\cfg}$.

We say that a grammar $\bnf$ is a \emph{valid} grammar, and that $valid(\bnf)$ evaluates to true, if it has a correct BNF syntax, and if $R$ is as defined above, and if all nonterminal symbols have at least one rule in their corresponding rule set.
% $r_i$, i.e., $1 \leq |r_i|\,\,. \forall 1 \leq i \leq n$.

% $Y$ consists of one or more sequences of terminal or non-terminal symbols, and the vertical bar ``$|$" separates different choices $y_i$ in $Y$.

% \subsection{Grammar Edge Set}

% The derivation space of a CFG can be represented as a \textit{grammar tree}~\cite{LLM-SYGUS}, which is the directed graph of sentential forms defined inductively as follows: the root of the tree is $S$, and for all $x, y \in (V \cup \Sigma)^*$, if $x \Rightarrow y$ in one derivation step with $x$ being a node of the tree, then $y$ is a child node of $x$.

% A \emph{grammar edge set} captures the immediate derivations from all non-terminals in the grammar tree. Given a CFG $\mathcal{G} = (V, \Sigma, R, S)$, we define the \emph{grammar edge set} as:
% $$
% \mathcal{E}^\mathcal{G} = \{e_1, e_2, \dots, e_\ell\},
% $$
% where $\ell = |V|$ is the number of non-terminals in $\mathcal{G}$, and each element $e_i$ represents the derivation information for a non-terminal.

% Each element $e_i$ in $\mathcal{E}^\mathcal{G}$, referred to as a \emph{grammar edge element}, is defined as  
% $$e_i = (x_i, Y_i).$$
% The component $x_i$ is a non-terminal in the grammar tree $T^\mathcal{G}$, satisfying $x_i \in V$. The set $Y_i$ consists of all child nodes of $x_i$, given by $Y_i = \{ y_{i,1}, y_{i,2}, \dots, y_{i,n_{x_i}} \},
% $ where each $y_{i,j}$ is derived from $x_i$ in a single step, i.e., $x_i \Rightarrow y_{i,j},\forall j \in [1, n_{x_i}]$. The number of direct derivations from $x_i$ is denoted as $n_{x_i}$, which defines the cardinality of $Y_i$.

% A \emph{grammar edge sequence} is a structural unit that captures all immediate derivations from all non-terminals in the grammar tree. Given a $\mathcal{G} = (V, \Sigma, R, S)$, we define a \emph{grammar edge compound} as $e = (x, Y)$, where $x$ being a node of the grammar tree $T^\mathcal{G}$, and $Y$ is the set of all child nodes of $x$.
% A \emph{grammar edge set} is 
% \begin{align*}
% \mathcal{E}{}^\mathcal{G} = \{ &e_i = (x_i, Y_i) \mid x_i \in V, \\ 
% & Y_i = \{ y_{i,1}, y_{i,2}, \dots, y_{i,n_{x_i}} \}, \\
% & x_i \Rightarrow y_{i,j}, 1 \leq j \leq n_{x_i} \},
% \end{align*}

% which represents all possible sentential forms that can be derived from $x_i$ in a single derivation step, and $n_{x_i}$ denotes the number of direct derivations from $x_i$. To provide an ordered structure, we organize the grammar edge set as a sequence, denoted as \emph{grammar edge sequence}: 
% $$\mathcal{E}^\mathcal{G} = \{e_1, e_2, \dots, e_\ell\},$$ 
% where $\ell = |V|$ is the number of non-terminals in $\mathcal{G}$. 



\subsection{Grammar Inference}
Grammar inference aims to learn grammar automatically from a set of examples~\cite{horning1969study,de2010grammatical,Stevenson_Cordy_2014,D_Ulizia_Ferri_Grifoni_2011}. In this paper, we focus on inferring a CFG, in BNF notation, from a small set of positive and negative examples. 

Given a set of positive examples $\mathcal{P}$ and negative examples $\mathcal{N}$, consisting of strings that must be, respectively, accepted and rejected, the objective is to infer a BNF grammar $\bnf$. The generated $\bnf$ should satisfy $\mathcal{P} \subseteq \languageOf{\bnf}$ and $\mathcal{N} \cap \languageOf{\bnf} = \emptyset$, which ensures $\bnf$ accepts all positive examples and rejects all negative examples.

% \subsection{Genetic Algorithm}
% Genetic Algorithm (GA) is a population-based metaheuristic method inspired by the principles of natural selection and genetics in biological evolution~\cite{Michalewicz_Schoenauer_1996,Katoch_Chauhan_Kumar_2021,Gendreau_Potvin_2019}. Typically, a GA relies on a set of genetic operators consisting of encoding, fitness, selection, crossover, and mutation~\cite{Katoch_Chauhan_Kumar_2021}. Based on these operators, a GA operates on a population of candidate solutions, refining the population over multiple generations, to find an optimal or near-optimal solution for a given problem.